\documentclass[10pt,a4paper]{amsart}
\usepackage{amsmath,amssymb,amsthm,graphicx}
\usepackage{amsmath}
\usepackage{amssymb}
\usepackage{hyperref}
\renewcommand{\theenumi}{\Alph{enumi}}
\renewcommand{\thesubsection}{\thesection\Alph{subsection}}

%%%%%%%%%%%%%%%%%%%%%%%%%%%%%%%%%%%%%%%%%%
% ENVIRONMENTS
%%%%%%%%%%%%%%%%%%%%%%%%%%%%%%%%%%%%%%%%%%%
\newtheorem{theorem}{Theorem}[section]
\newtheorem{prop}[theorem]{Proposition}
\newtheorem{lemma}[theorem]{Lemma}
\newtheorem{remark}[theorem]{Remark}
\newtheorem{corollary}[theorem]{Corollary}
\newtheorem{definition}[theorem]{Definition}
\newtheorem{example}[theorem]{Example}
%
\def\beq{\begin{equation}}
\def\eeq{\end{equation}}
\def\be{\begin{equation*}}
\def\ee{\end{equation*}}
%%%%%%%%%%%%%%%%%%%%%%%%%
% Beginning of the text %
%%%%%%%%%%%%%%%%%%%%%%%%%
%\begin{document}



%\documentclass[10pt]{article}
\usepackage{amssymb}
\usepackage{amsfonts}
%\topmargin 0.3in
%\textwidth 6in
%\textheight 9in%\baselineskip
%\hoffset=-50pt
%\voffset=-2pt
%\newenvironment{proof}[1][Proof]
%{\ \rule{0.5em}{0.5em}}


\title{Dedekind-finite Cardinals having countable partitions}
\author{Supakun Panasawatwong and J K Truss} 
\date{This paper is an amplified version of part of the first author's PhD thesis at the University of Leeds}%\footnotemark}
\title{Every omega-tree with finite levels contains a leafless balanced subtree with omega levels}
\author{Supakun Panasawatwong, J. K. Truss}
\date{arxiv:2204.05700v1; CC BY 4.0}
\begin{document}
\maketitle
\noindent\textit{Editorial note.} All delivered mathematical text is the author's own, copied from the pinned publisher source sources/dedekind-finite-trees-2204.05700v1/source0.tex; no statement or proof step was rewritten, reordered or removed.\par\medskip
\noindent\textit{Editorial note.} Preamble provenance: the retained preamble is the source's own layout adapter, if any, plus the author's macro block at lines 5-44 of sources/dedekind-finite-trees-2204.05700v1/source0.tex (everything between the publisher class metadata and \textbackslash{}begin\{document\}, comments excluded). Nothing was added, no macro was invented, and the publisher document class is deliberately not reproduced; the excerpt serves only as preamble, never as a span coordinate.\par\medskip
\noindent\textit{Editorial note.} External references are navigation pointers only; every deliberately kept step of the author's argument stays inside the counted spans. The displayed numbers of those pointers are the author's own numbering under the author's \textbackslash{}newtheorem declarations (the source has no subsections, so numbering is section.index), computed mechanically by the lane's scanner.\par\medskip
\noindent\textit{Editorial note.} The counted proof contains no reference to any labelled item and no bibliography citation, so the card needs neither a navigation stub nor a bibliography excerpt.\par\medskip
\noindent\textit{Editorial note.} Source metadata: this directory stages no SOURCE.json; the title and author list were taken from the publisher abstract page kept beside it, and the card points at a per-card licence reuse record in its own task directory because the staged licence evidence carries no status field.\par\medskip
\noindent\textit{Allowed standard prerequisites (editorial).}\begin{itemize}
\item The notion of a balanced subtree and of the degree of a node on a level. Defined in the paragraphs preceding the counted statement (the discussion ending at source line 489). Fixes what has to be produced and what the pruning at each stage must achieve.
\item A sequence of natural numbers in which every number occurs infinitely often. Used inside the counted proof at lines 494-495 and 508. Supplies the schedule according to which the successive prunings are performed and why each level is handled infinitely often.
\end{itemize}\medskip
\begin{lemma} Any tree $(T, \le)$ with $\omega$ levels, all finite, has a balanced subtree, having no leaves,
and also having $\omega$ levels. \label{4.1}  \end{lemma}

\begin{proof} Let $S$ be a sequence of natural numbers such that every number occurs infinitely often, 
say $S = \langle k_n: n \in \omega \rangle$. We construct a decreasing sequence $T_n$ of subtrees 
of $T$, each having $\omega$ levels, such that for each $n$, $T_n$ is pruned on level $k_n$ so that 
every member has the same degree (number of immediate successors) on that level. Let $L_n$ be the 
$n$th level of $T$.

First let $T_0 = T$. Now suppose that $T_n$ has been constructed. Let $X_n$ be the set of
members $x$ of height $k_n$ in $T_n$ such that $\{y \in T_n: x \le y\}$ is infinite, and subject to 
this have minimum degree (since $T_n$ has $\omega$ levels, at least one node on that level has 
infinitely many points above it). Let $T_{n+1}$ be the subtree obtained by removing all members 
of $T_n$ $\ge$ some member of $L_{k_n} \setminus X_n$. Then $T_{n+1}$ is a subtree of $T_n$ also having
$\omega$ levels. Let $T^* = \bigcap_{n \in \omega} T_n$. It remains to show that $T^*$ is balanced 
and has no leaves.

Consider any $n$. Since each member of $\omega$ is listed infinitely often, $\{m: k_m = n\}$ is
infinite. Let $L_n^m$ be the $n$th level of $T_m$. Then $L_n^0 \supseteq L_n^1 \supseteq L_n^2 \ldots$
is a decreasing sequence of non-empty finite sets, so is eventually constant, and clearly this equals the 
$n$th level of $T^*$. Furthermore every $x$ in the $n$th level of $T^*$ has at least one successor in
every $T_m$, and hence also in $T^*$, so is not a leaf, and for the same reason, its degree is 
constant among nodes on that level. Hence $T^*$ is a balanced subtree of $T$.
  \end{proof} 

\end{document}
